%%%%%%%%%%%%%%%%%%%%%%%%%%%%%%%%%%%%%%%%%%%%%%%%%%%%%%%%%%%%%%%%%%%%%%%%%%%%%
%% Chapter 1 — Introduction
%% Budget: 3 minutes / 3 slides
%%   S1  Motivation         — application class + heterogeneous-tool reality
%%   S2  Research gap       — capability table + central RQ + sub-RQ keywords
%%   S3  Contribution       — what was built + forward pointers
%%%%%%%%%%%%%%%%%%%%%%%%%%%%%%%%%%%%%%%%%%%%%%%%%%%%%%%%%%%%%%%%%%%%%%%%%%%%%

\section{Introduction}

%---------------------------------------------------------------------------
\begin{frame}[fragile]{Heterogeneous Device--Human Systems}
    \centering
    \vspace{-0.35em}
    \begingroup
        \resizebox{!}{0.86\textheight}{%
            % Reusable TikZ body for the exoskeleton motivation figure.
% Define \exopath and \exologopath before input if assets are not resolved from
% this directory.
\providecommand{\exopath}{figures/introduction/}
\providecommand{\exologopath}{assets/logos/}

\begin{tikzpicture}[pres]
    \coordinate (min) at (0, 0);
    \coordinate (max) at (16, 10);
    \useasboundingbox (0,0.4) rectangle (16, 9.6);

    % Grid help lines (optional, for alignment)
    % \draw[help lines, step=0.5, draw=gray!50, very thin] (min) grid (max);

    \begin{scope}[yshift=0cm]
        % Figure content
        \node[inner sep=0, anchor=south west] (img) at (4.5,1.5) {%
            \includegraphics[height=7.5cm]{\exopath exo.png}%
        };
        
        % Ground line and hatch
        \coordinate (A) at (4.5, 1.5);
        \coordinate (B) at (11.75, 1.5);
        \draw[pres ground line] (A) -- (B);
        \fill[pres ground hatch] ([yshift=-5mm]A) rectangle (B);
        
        % Biomechanics subsystem box
        \coordinate (muscleLabel) at (2, 7.25);
        \node[pres subsystem box] (n-bio) at (muscleLabel) {%
            \begin{tabular}{@{}c@{}}
                \textbf{Human}\\
                \textbf{Biomechanics}\\[1pt]
                {\footnotesize\itshape Musculoskeletal}\\
                {\footnotesize\itshape Modeling}\\[2pt]
                \IfFileExists{\exologopath OpenSim.png}
                    {\includegraphics[height=10mm]{\exologopath OpenSim.png}}
                    {{\scriptsize OpenSim}}
            \end{tabular}%
        };
        
        % Control subsystem box
        \coordinate (controlLabel) at (14, 5);
        \node[pres subsystem box] (n-control) at (controlLabel) {%
            \begin{tabular}{@{}c@{}}
                \textbf{Actuators, Control,}\\
                \textbf{and Sensors}\\[1pt]
                {\footnotesize\itshape Equation-Based}\\
                {\footnotesize\itshape Modeling}\\[2pt]
                \IfFileExists{\exologopath Modelica.pdf}
                    {\includegraphics[height=10mm]{\exologopath Modelica.pdf}}
                    {{\scriptsize Modelica}}
            \end{tabular}%
        };
        
        % FEM subsystem box
        \coordinate (structureLabel) at (2, 2.75);
        \node[pres subsystem box] (n-structure) at (structureLabel) {%
            \begin{tabular}{@{}c@{}}
                \textbf{Mechanical}\\
                \textbf{Structure}\\[1pt]
                {\footnotesize\itshape Finite Element}\\
                {\footnotesize\itshape Analysis}\\[2pt]
                \IfFileExists{\exologopath NGSolve.png}
                    {\includegraphics[height=8mm]{\exologopath NGSolve.png}}
                    {{\scriptsize NGSolve}}\\
                {\scriptsize\itshape Netgen/NGSolve}
            \end{tabular}%
        };
        
        % Leaders from boxes to image
        \coordinate (p-muscle) at (8.1, 6.8);
        \coordinate (p-control) at (9, 7);
        \coordinate (p-structure) at (10, 1.5);

        \draw[pres leader] (n-bio.east) -- (p-muscle);
        \draw[pres leader] (n-control.west) -- (p-control);
        \draw[pres leader] (n-structure.east) -- (p-structure);

        \draw[pres coupling] (n-bio.north) --++ (0, 0.5) --++ (12, 0) -- (n-control.north);
        \draw[pres coupling] (n-structure.south) --++ (0, -0.5) --++ (12, 0) -- (n-control.south);
        \draw[pres coupling] (n-structure.north) -- (n-bio.south);
    \end{scope}
\end{tikzpicture}
%
        }
    \endgroup
\end{frame}

%---------------------------------------------------------------------------
\begin{frame}{Co-Simulation Challenges}
  \begin{columns}[T]
    \begin{column}{0.55\textwidth}
      \centering
      \resizebox{\linewidth}{!}{%
        \begin{tikzpicture}[pres]
  % Force identical bounding box across all four panels
  \coordinate (min) at (0,0.7);
  \coordinate (max) at (6.8, 7.3);
  \useasboundingbox (min) rectangle (max);
  % \draw[help lines, step=0.5] (min) grid (max);

  % SCC Block
  \draw[pres loop area] (3.2, 0.7) rectangle (6.8, 7.3);
  \node[pres loop label] at (5.15, 4) {algebraic loop};

  % Component A
  \begin{scope}[shift={(0,4.5)}]
    \draw[pres orchestration component] (0,0) rectangle (2,2.5);
    \node[anchor=north, font=\Large\bfseries] (A) at (1,2.45) {A};
    \node[anchor=north, inner sep=1pt, font=\normalsize] at (A.south) {\faCog};
    \node[pres port] (A-y1) at (1.75, 1.75) {};
  \end{scope}

  % Component B
  \begin{scope}[shift={(4,4.5)}]
    \draw[pres orchestration loop component] (0,0) rectangle (2,2.5);
    \node[anchor=north, font=\Large\bfseries] (B) at (1,2.45) {B};
    \node[anchor=north, inner sep=1pt, font=\normalsize] at (B.south) {\faCog};
    \node[pres port] (B-u1) at (0.25,1.75) {};
    \node[pres port] (B-u2) at (0.25,0.75) {};
    \node[pres port] (B-y2) at (1.75,0.75) {};
  \end{scope}

  % Component C
  \begin{scope}[shift={(4,1)}]
    \draw[pres orchestration loop component] (0,0) rectangle (2,2.5);
    \node[anchor=north, font=\Large\bfseries] (C) at (1,2.45) {C};
    \node[anchor=north, inner sep=1pt, font=\normalsize] at (C.south) {\faCog};
    \node[pres port] (C-y1) at (0.25,1.75) {};
    \node[pres event port] (C-y2) at (0.25,0.75) {};
    \node[pres port] (C-u1) at (1.75,1.75) {};
  \end{scope}

  % Component D
  \begin{scope}[shift={(0,1)}]
    \draw[pres orchestration component] (0,0) rectangle (2,2.5);
    \node[anchor=north, font=\Large\bfseries] (D) at (1,2.45) {D};
    \node[anchor=north, inner sep=1pt, font=\normalsize] at (D.south) {\faCog};
    \node[pres event port] (D-listener) at (1.75,0.75) {};
  \end{scope}

  % Connections
  \draw[pres signal] (A-y1) -- (B-u1);
  \draw[pres signal] (B-y2) --++(0.25+0.5,0) --++(0,-2.5) --(C-u1);
  \draw[pres signal] (C-y1) --++(-0.25-0.5,0) --++(0,2.5) -- (B-u2);
  \draw[pres event signal] (C-y2.west) -- (D-listener.east);

  % Badges
  \node[pres badge] (b1) at (2.8,6) {1};

  \node[pres badge] (b2) at (3.9,4) {2};

  \node[pres badge] (b3) at (2.8,2) {3};

\end{tikzpicture}
%
      }
    \end{column}
    \begin{column}{0.43\textwidth}
        \begin{enumerate}
            \item \textbf{Execution order}\\
            Which simulation unit runs first?
            \vspace{0.8em}
            \item \textbf{Algebraic loops}\\
            Cyclic dependencies require an interface solve.
            \vspace{0.8em}
            \item \textbf{Hybrid events}\\
            Detect, localize, and restore consistent state.
        \end{enumerate}
    \end{column}
  \end{columns}
\end{frame}



%---------------------------------------------------------------------------
% \begin{frame}{Research Gap}
%     \begin{itemize}
%         \item \ac{FMI}-based frameworks handle \ac{FMU} coupling and algebraic loops.
%         \item None combines native \ac{FMU}, OpenSim, and \ac{FEM} backends with hybrid event handling and runtime model switching.
%     \end{itemize}

%     \vspace{0.3em}
%     \begin{table}[h]
%         \centering\footnotesize
%         \begin{tabular}{l c c c c c c}
%             \toprule
%             \textbf{Framework}  & \textbf{\ac{FMU}}    & \textbf{OpenSim} & \textbf{\ac{FEM}}
%                                 & \textbf{Alg.\ loops} & \textbf{Hybrid}  & \textbf{Switching}                                            \\
%             \midrule
%             OMSimulator         & \checkmark           & ---              & ---                & \checkmark & ($\checkmark$) & ---        \\
%             INTO-CPS Maestro    & \checkmark           & ---              & ---                & \checkmark & ($\checkmark$) & ---        \\
%             CoFMPy              & \checkmark           & ($\checkmark$)   & ($\checkmark$)     & \checkmark & ---            & ---        \\
%             \textbf{\syssimx{}} & \checkmark           & \checkmark       & \checkmark         & \checkmark & \checkmark     & \checkmark \\
%             \bottomrule
%         \end{tabular}
%     \end{table}
% \end{frame}

%---------------------------------------------------------------------------
% \begin{frame}{Contribution}
%     \begin{block}{\syssimx{} --- a Python framework for heterogeneous hybrid co-simulation}
%         This thesis designs, implements, and evaluates:
%         \begin{itemize}
%             \item a unified component abstraction for heterogeneous subsystem models,
%             \item graph-based structural analysis (dependency detection, execution order, algebraic loops),
%             \item hybrid event handling (detection, localization, dispatch, cascaded events),
%             \item a multi-model component for runtime switching with state synchronization.
%         \end{itemize}
%     \end{block}

%     \begin{itemize}
%         \item Evaluated on a controlled-pendulum case study combining \acp{FMU}, OpenSim, and \ac{FEM} models in one closed-loop scenario.
%         \item Evidence path: architecture $\to$ implementation $\to$ case study.
%     \end{itemize}
% \end{frame}
